\documentclass[11pt]{article}
\usepackage[a4paper,margin=27mm]{geometry}
\usepackage[T1]{fontenc}
\usepackage{lmodern,microtype,amsmath,amssymb,amsthm,mathtools,enumitem}
\usepackage[hidelinks,hypertexnames=false]{hyperref}
\newtheorem{theorem}{Theorem}
\newtheorem{lemma}[theorem]{Lemma}
\newtheorem{proposition}[theorem]{Proposition}
\newtheorem{corollary}[theorem]{Corollary}
\theoremstyle{remark}\newtheorem{remark}[theorem]{Remark}
\newcommand{\R}{\mathbb R}
\newcommand{\Dscr}{\mathcal D}
\newcommand{\Escr}{\mathcal E}
\newcommand{\supp}{\operatorname{supp}}
\newcommand{\one}{\mathbf 1}
\title{Closed Choice Beyond Ordinary Integrals}
\author{}\date{9 October 2026}
\begin{document}\maketitle

\noindent
The unrestricted problem remains open. This note strengthens the obstruction
in \emph{Closed Choice for Independent Rankings} (CC), and shows why one
apparently promising proof strategy cannot work.
\begin{itemize}[leftmargin=*,itemsep=3pt]
\item Finitely additive penalties do not escape the obstruction. Even an
incomplete ranking with a strictly guidance-respecting finitely additive
weak extension on regular models cannot leave the vacuous model undominated
when $n\ge3$ (Theorem~\ref{thm:fa}).
\item Nevertheless, (S) does \emph{not} force an ordinary additive weak
extension on regular models. A complete, real-valued ranking satisfies
(G), (IB), (LA), (S), and (C), but has no such extension
(Proposition~\ref{prop:example}). Thus CC's Theorem~1 cannot become a general
impossibility proof simply by deriving its extra hypothesis from (S).
\item Closed choice forces continuity of any measure-continuous real primary
along two-world edges. It also rules out nontrivial priorities based solely
on advice in arbitrarily small sets (Section~\ref{sec:primary}).
\end{itemize}
The first result uses the additive obstruction already proved in v3.
It does not assume that independence entails an additive representation.

\section{Setting and the facts needed from the earlier notes}

There are $n$ worlds. Put
$\Delta=\{p\in\R^n:p\ge0,\ \sum_i p_i=1\}$ and
$H=\{g\in\R^n:g_1>0\}$. We normalize Gaussian measure on $H$ to have mass
one and call it $\mu$. All models are measurable, and equality means equality
outside a null set. At each binary problem $\{g,-g\}$ a model accepts $g$
only, accepts neither, accepts $-g$ only, or accepts both. Its guidance rank
$r_i(g,\Dscr)$ at world $i$ is respectively $0,1,2,3$ when the advice is
good, silent, bad, or contradictory. Which decisive recommendation is good
depends on the sign of $g_i$; coordinate hyperplanes are null.

The ideal is $D_i=\{g:g_i>0\}$. A model is \emph{regular} if it never
contradicts itself and is correct on unanimous problems. Every coherent
model is regular. Write $H'$ for the mixed problems in $H$.

A coherent model has, up to a null set, the form
\[
 \Dscr_M=\{g:\underline{P_M}(g)>0\}\cup\{g\ge0:g\ne0\},
 \qquad \underline{P_M}(g)=\min_{m\in M}m\cdot g,
\]
where $M$ is a nonempty compact convex subset of $\Delta$.
The vacuous model is $\Dscr_\Delta$.

At each world $\preceq_i$ is a preorder, with smaller meaning better.
\textbf{(G)} preserves weak pointwise improvement in guidance rank, and makes
the comparison strict if the improvement occurs on a positive-measure set.
\textbf{(IB)} says that replacing advice shared by two models on any measurable
set of binary problems leaves their comparison unchanged.

A model $C$ is competitive on $S$ if, for every $\Dscr$, some $i\in S$
satisfies $C\preceq_i\Dscr$. Requirement \textbf{(S)} supplies coherent
undominated models $C_p$, $p\in\Delta$, competitive on $\supp p$, with
\[
 A_i(\Dscr)=\{p:C_p\preceq_i\Dscr\}\quad\text{closed for every }i,\Dscr.
 \tag{S}
\]
Requirement \textbf{(LA)} asks that coherent undominated models exist.
Requirement \textbf{(C)} says that a model tied everywhere with a coherent
undominated model is itself coherent and undominated.

We need the following consequence of (S), which does not require complete
rankings.
\begin{lemma}\label{lem:dominator}
Every model has a selected coherent weak dominator at all worlds. Every
undominated model is tied everywhere with a selected model, and is therefore
competitive on some nonempty support.
\end{lemma}
\begin{proof}
Fix $\Dscr$. The closed sets $A_i=A_i(\Dscr)$ cover each face by its own
indices, by competitiveness. Put $d_i(p)=\operatorname{dist}(p,A_i)$ and
\[
 F_i(p)=\frac{p_i+d_i(p)}{1+\sum_jd_j(p)}.
\]
Each $A_i$ contains $e_i$. A fixed point exists by Brouwer's theorem.
At a fixed point, $d_i=p_i\sum_jd_j$. Some supported $i$ belongs to $A_i$,
so $d_i=0$ and $p_i>0$, forcing every $d_j=0$. This gives a common point of
the $A_i$. If $\Dscr$ is undominated, its weak dominator can be strictly
better nowhere, so all comparisons are ties. Transitivity transfers
competitiveness across those ties.
\end{proof}

We also use the following established result from v3, Appendix C:
\begin{quote}
For ordinary integral penalties that strictly order good, silence, bad,
and contradiction and have an integrable envelope, suppose a coherent
expected-loss minimizer exists at every prior. If $n\ge3$, no
$\Dscr_M$ with $M$ containing three simplex vertices is an expected-loss
minimizer at any prior. Each such model is strictly dominated at all worlds
for these penalties.
\end{quote}
The last conclusion follows by separation: the closure of the attainable
loss vectors is convex, by atomless splitting. Absence of a strict
dominator gives a supporting nonnegative weight vector, hence an
expected-loss minimum. This is the only earlier obstruction used below.

\section{Finitely additive penalties still cannot give full coverage}

Finite additivity allows an evaluation to give positive weight to sets
whose Gaussian measures tend to zero. Thus these penalties need not be
continuous in measure. The next theorem includes this possibility.

Let $\nu_{ik}$, $i\le n$, $1\le k\le3$, be finite nonnegative finitely
additive measures on the measurable subsets of $H$, satisfying
\begin{equation}
 \nu_{ik}(A)=0\text{ if }\mu(A)=0,
 \qquad \nu_{ik}(A)>0\text{ if }\mu(A)>0.
 \label{eq:positive}
\end{equation}
Define the finite losses
\begin{equation}
 L_i(\Dscr)=\sum_{k=1}^3\nu_{ik}\{g:r_i(g,\Dscr)\ge k\}.
 \label{eq:fa}
\end{equation}
The three measures charge the successive deteriorations in guidance.
Ordinary local penalties are the special case where the measures have
integrable densities. Equation~\eqref{eq:fa} satisfies (G) and (IB).

\begin{theorem}\label{thm:fa}
Let $n\ge3$. Suppose a ranking satisfies (G), (LA), and (S), and for each
world it has a loss of the form \eqref{eq:fa} such that
\begin{equation}
 \Dscr\preceq_i\Escr\ \Longrightarrow\ L_i(\Dscr)\le L_i(\Escr)
 \qquad\text{for regular }\Dscr,\Escr.
 \label{eq:weak}
\end{equation}
Then every $\Dscr_M$ whose credal set contains three simplex vertices is
dominated by a coherent model. If the ranking is represented by $L$ itself,
the coherent dominator is strictly better at every world.
\end{theorem}

The proof separates a penalty into an ordinary integral and a part that
can be held fixed while advice is changed almost everywhere else. Two
short lemmas make this precise.

\begin{lemma}\label{lem:split}
There are strictly guidance-ordered ordinary integral losses $L_i^c$ and
nonnegative finitely additive losses $L_i^s$ with $L_i=L_i^c+L_i^s$.
There are measurable $Z_m$ with $\mu(Z_m)\to0$ such that, if $X$ gives the
advice of $A$ on $Z_m$ and of $B$ elsewhere, then
\[
 L_i^s(X)=L_i^s(A),\qquad
 |L_i^c(X)-L_i^c(B)|\le\int_{Z_m}h_i\,d\mu\longrightarrow0,
\]
where $h_i$ is a fixed integrable function. The bound is uniform in $A,B$.
\end{lemma}
\begin{proof}
The Yosida--Hewitt decomposition and its concentration property
\cite[Theorem 1 and its proof]{PP} give
\[
 \nu_{ik}(A)=\int_A f_{ik}\,d\mu+\pi_{ik}(A),\qquad f_{ik}\ge0,
\]
and sets $Z_m$ with $\mu(Z_m)\to0$ and
$\pi_{ik}(H\setminus Z_m)=0$ for every $i,k$. A finite union gives common
sets for all the measures.

In fact $f_{ik}>0$ a.e. If it vanished on a positive-measure $A$, choose
$m$ with $\mu(Z_m)<\mu(A)$. Then $B=A\setminus Z_m$ has positive measure,
but both terms in $\nu_{ik}(B)$ vanish, contradicting
\eqref{eq:positive}. Summing the densities over the exceeded ranks gives
$L_i^c$, with strictly positive successive increments and envelope
$h_i=\sum_k f_{ik}$. Summing the pure parts gives $L_i^s$.
Concentration on $Z_m$ proves the first displayed equality; the envelope
proves the bound.
\end{proof}

\begin{lemma}\label{lem:compact}
The coherent models are sequentially compact for convergence in measure
of their advice. Each $L_i^c$ is continuous for this convergence.
\end{lemma}
\begin{proof}
The nonempty compact convex subsets of the compact simplex are compact
in Hausdorff distance. If $M_m\to M$, then
$\underline{P_{M_m}}(g)\to\underline{P_M}(g)$ for every $g$.
The zero set of $\underline{P_M}$ is null: on each line parallel to
$\one=(1,\ldots,1)$,
$\underline{P_M}(g+t\one)=\underline{P_M}(g)+t$, so there is just one zero;
apply Fubini. The same applies to $-g$. Advice therefore converges a.e.,
and hence in measure. Finally,
$|L_i^c(A)-L_i^c(B)|\le\int_{\{a_A\ne a_B\}}h_i\,d\mu$.
\end{proof}

\begin{proof}[Proof of Theorem~\ref{thm:fa}]
\emph{First obtain coherent expected-loss minima for $L^c$.}
Fix a prior $p$. Among regular models let
\[
 a=\min_B\sum_i p_i L_i^c(B),\qquad
 b=\inf_A\sum_i p_i L_i^s(A).
\]
The first minimum exists by choosing a least costly advice at each problem.
Its minimizer $B$ is regular. All quantities are finite and nonnegative.
Choose regular $A_m$ with pure expected loss tending to $b$, and let $X_m$
use $A_m$ on $Z_m$ and $B$ elsewhere. Then
\[
 \sum_i p_i L_i(X_m)\longrightarrow a+b.
\]
Lemma~\ref{lem:dominator} supplies a selected coherent $C_m$ with
$C_m\preceq_i X_m$ for every $i$. Both are regular, so \eqref{eq:weak} gives
\[
 a\le\sum_i p_iL_i^c(C_m)
 \le\sum_i p_iL_i(X_m)-b\longrightarrow a.
\]
By Lemma~\ref{lem:compact}, a subsequence converges in measure to a coherent
model attaining $a$. Thus coherent $L^c$-Bayes models exist at every prior.
Notice that we have not assumed that a finitely additive expected-loss
infimum is attained.

\emph{Next transfer strict domination back to $L$.}
The v3 obstruction applies to $L^c$. For any target $T=\Dscr_M$ as in the
theorem, there is a regular $B$ with $L_i^c(B)<L_i^c(T)$ at every world.
Let $X_m$ agree with $T$ on $Z_m$ and with $B$ elsewhere. Then
\[
 L_i(X_m)-L_i(T)=L_i^c(X_m)-L_i^c(T)
 \longrightarrow L_i^c(B)-L_i^c(T)<0.
\]
For large $m$, $X_m$ strictly improves every $L_i$.

If $T$ were undominated for the original ranking, Lemma~\ref{lem:dominator}
would make it competitive on some support. In particular
$T\preceq_i X_m$ for some $i$, contrary to \eqref{eq:weak}. So $T$ is
dominated. Applying Lemma~\ref{lem:dominator} to a dominator gives a coherent
dominator by transitivity. When $L$ represents the ranking, $X_m$ itself
strictly dominates $T$ at every world; a selected weak dominator of $X_m$
retains all these strict inequalities.
\end{proof}

\section{Closed choice does not force an ordinary additive extension}

The preceding theorem is strictly stronger than the ordinary-integral
obstruction. The following example satisfies closed choice but lies outside
the older theorem's extra hypothesis, even on regular models.

Write $g=R\theta$ under normalized Gaussian measure on $H$. Direction
$\theta$ and radius $R$ are independent. If $F$ is the continuous distribution
function of $R$, then $t=F(R)$ is uniform on $(0,1)$ and independent of
$\theta$. Fix a nonprincipal ultrafilter $\mathcal U$ on the positive
integers and define, for measurable $A\subseteq H$,
\[
 \eta(A)=\lim_{m\to\mathcal U}m\,\mu(A\cap\{t<1/m\}),
 \qquad \nu=\tfrac12\mu+\tfrac12\eta.
\]
An ultrafilter limit chooses a limit consistently for every bounded real
sequence; it preserves sums, inequalities, and ordinary limits. Thus $\eta$
is a finitely additive probability, vanishes on $\mu$-null sets, and gives
mass one to every $\{t<1/m\}$. It is not countably additive. On functions
of direction alone its integral agrees with $\mu$.

Put $b_i(g)=|g_i|/\|g\|$. At world $i$ assign the four ranks the costs
\begin{equation}
 0,\quad \tfrac34 b_i(g),\quad b_i(g),\quad 2b_i(g),
 \qquad L_i(\Dscr)=\int_H\ell_i(g,a_\Dscr(g))\,d\nu.
 \label{eq:example}
\end{equation}
These are finite losses of the form \eqref{eq:fa}.

\begin{proposition}\label{prop:example}
For every $n\ge2$, the complete ranking represented by \eqref{eq:example}
satisfies (G), (IB), (LA), (S), and (C), with
$C_p=\Dscr_{\{p\}}$. It has no ordinary strictly guidance-ordered additive
weak extension on regular models, at any world.
\end{proposition}
\begin{proof}
Strict (G) follows from the $\mu/2$ term, and (IB) from finite additivity.
For a prior $p$, let $B_+$ and $B_-$ be the expected costs of recommending
$g$ and $-g$. The cost of silence is
$3(B_++B_-)/4$, strictly greater than $\min(B_+,B_-)$ except on a null set.
Contradiction is worse still. The preferred decisive advice is determined
by the sign of $p\cdot g$. Its tie hyperplane is null for both $\mu$ and
$\eta$, since it depends only on direction. Thus $C_p$ is the unique Bayes
model modulo null sets: any different model has a strictly positive
expected-loss gap under the $\mu/2$ term.

Bayes optimality gives competitiveness on $\supp p$. Uniqueness gives
undominatedness, including for boundary priors. For coherent models
both the advice and the costs depend only on direction, so
\[
 L_i(C_p)=\int_H\ell_i(g,a_{C_p}(g))\,d\mu.
\]
This is continuous in $p$ by dominated convergence. Therefore all sections
$\{p:L_i(C_p)\le L_i(\Dscr)\}$ are closed, proving (S).
Lemma~\ref{lem:dominator} now makes every undominated model tied everywhere
with a selected $C_p$; uniqueness of that $p$-Bayes model makes it equal to
$C_p$ a.e. The same reasoning applies to any model tied with it, proving (C).

Fix a world $i$ and a mixed cone $Q\subseteq H'$ of positive measure.
Let $B_m$ be regular, silent on $Q\cap\{t<1/m\}$, and ideal at $i$
elsewhere. Then $B_m\to D_i$ in measure but
\[
 L_i(B_m)\ge a\coloneqq\frac38\int_Q b_i\,d\mu>0.
\]
Choose a regular $E\ne D_i$ that is silent on a sufficiently small
positive-measure subset of $H'\cap\{t>1/2\}$ and ideal at $i$ elsewhere.
Then $0<L_i(E)<a$, so $E\prec_i B_m$ for every $m$.
An ordinary integral weak extension $\Lambda_i$ on regular models would
give
\[
 \Lambda_i(E)\le\Lambda_i(B_m)\longrightarrow\Lambda_i(D_i),
\]
contradicting its strict reward for the guidance improvement from $E$ to
$D_i$. This proves the claimed failure of an ordinary weak extension.
\end{proof}

This is not a coverage construction: Theorem~\ref{thm:fa} excludes the
vacuous target for $n\ge3$. It shows, however, that neither completeness,
real-valuedness, independence, nor closed choice supplies the ordinary
extension assumed by CC's Theorem~1. The discontinuity can occur on mixed
problems themselves, not just on errors in unanimous problems.

\section{What closed choice forces on real primary criteria}\label{sec:primary}

A \emph{primary} at world $i$ is a real function $\Phi$ such that
$\Phi(A)<\Phi(B)$ implies $A\prec_i B$. It automatically weakly respects
$\preceq_i$: otherwise a weak comparison would contradict a strict reverse
comparison. The remaining ranking may be incomplete.

The next lemma uses only small guidance improvements and the closed
sections in (S). In particular, it does not assume (IB).

\begin{lemma}[Continuity along an edge]\label{lem:edge-continuity}
Assume (G) and (S), and let $\Phi$ be a primary at $i$, continuous in measure
on regular models. For $j\ne i$, the function
$p\mapsto\Phi(C_p)$ is continuous on the edge $[e_i,e_j]$.
\end{lemma}
\begin{proof}
First, competitiveness and strict (G) force every selected credal set on
this edge to lie on the edge. Otherwise correcting advice where both edge
worlds are positive and an unsupported world is sufficiently negative
strictly improves both supported worlds. The endpoint selections are
therefore $D_i,D_j$.

Write $f(p)=\Phi(C_p)$. It is lower semicontinuous. Indeed, if
$p_m\to p$ and $f(p_m)\le f(p)-\epsilon$, a regular model $E$ can be chosen
with $f(p)-\epsilon<\Phi(E)<f(p)$. To obtain $E$, continuously replace one
regular model by another on a growing measurable set and use continuity
of $\Phi$. Then $C_{p_m}\prec_i E$ while $E\prec_i C_p$, contradicting the
closed section $\{q:C_q\preceq_i E\}$.

For upper semicontinuity suppose instead $f(p_m)\ge f(p)+\epsilon$.
If $C_p=D_j$, this is impossible: at world $i$, $D_j$ is guidance-worst
among coherent edge models. Otherwise $C_p$ can be strictly improved at
$j$ on an arbitrarily small positive-measure set. Choose such a regular
$E\prec_j C_p$ with $\Phi(E)<f(p)+\epsilon$.
Then $C_{p_m}\not\preceq_i E$, by the primary property. Competitiveness
on the edge forces $C_{p_m}\preceq_j E$. Closedness gives
$C_p\preceq_j E$, a contradiction.
\end{proof}

\begin{remark}
The upper-semicontinuity argument extends to a face $\Delta_S$ at any
selected $C_p=\Dscr_K$ with $\max_{m\in K}m_i>0$.
There is then an open set of problems with an option positive at every world
in $S\setminus\{i\}$ that $C_p$ fails to recommend.
A correction on a small subset strictly improves all these worlds.
Competitiveness and a finite subsequence argument give the same contradiction.
Thus at a selected vacuous model every measure-continuous primary is
continuous along the simplex. This does not establish continuity of the
rankings on all models.
\end{remark}

There is an even sharper conclusion when a primary ignores some
positive-measure part of every region. This includes priorities determined
by purely finitely additive measures or by finitely many null-invariant ultrafilter
evaluations.

\begin{theorem}[No priority confined to arbitrarily small sets]\label{thm:blind}
Assume (G) and (S). Let $\Phi$ be a primary at world $i$ with this property:
for every positive-measure $Q\subseteq H'$ there is a positive-measure
$B\subseteq Q$ such that changing regular advice only on $B$ never changes
$\Phi$. Then $\Phi(C_p)$ is constant on every edge $[e_i,e_j]$.
In particular, $\Phi(D_i)=\Phi(D_j)$ for every $j$.
\end{theorem}
\begin{proof}
Fix an edge and put $f(p)=\Phi(C_p)$. If $C_p\ne D_i$, improve its advice
at $i$ on a positive-measure subset ignored by $\Phi$. The resulting regular
$E$ has $E\prec_i C_p$ and $\Phi(E)=f(p)$. Closedness excludes
$C_q\preceq_i E$ near $p$, forcing $f(q)\ge f(p)$. If $C_p=D_i$, the same
inequality holds globally because $D_i$ is guidance-best at $i$.

Similarly, if $C_p\ne D_j$, improve it at $j$ on an ignored subset,
obtaining $F\prec_j C_p$ with $\Phi(F)=f(p)$. Near $p$, closedness excludes
$C_q\preceq_j F$. Competitiveness forces $C_q\preceq_i F$, and therefore
$f(q)\le f(p)$. If $C_p=D_j$, this inequality holds globally because $D_j$
is guidance-worst at $i$ among coherent edge models.

Thus $f$ is locally constant everywhere on the connected edge, so it is
constant.
\end{proof}

For example, if finitely many pure charges determine a primary, choose
a common concentrating set $Z$ with $\mu(Z)<\mu(Q)$. Its complement inside
$Q$ is an ignored positive-measure set. Theorem~\ref{thm:blind} therefore
rules out the construction as soon as the primary distinguishes two edge
ideals. The later criteria and possible incomparabilities do not help.

\section{The remaining gap}

CC's conjecture has not been proved or refuted. The strongest conclusion
here is that a successful ranking cannot have even a finitely additive
strict-guidance weak extension on regular models. This includes incomplete
rankings; simply dropping countable additivity is insufficient.

There are two different obstacles to completing the argument.
First, (IB) cancels common advice, but it does not supply the finite
cancellation conditions needed for an additive weak extension. The
six-region example in BA already witnesses that distinction for partial
rankings. Second, closed choice does not repair it by forcing an ordinary
integral extension: Proposition~\ref{prop:example} explicitly disproves
that inference, even on regular models.

A general proof would still need to derive a suitable stronger consistency
property from competitive coherent choice, or directly contradict that
choice without an additive representation. No such implication is assumed
in the results above. A construction outside these classes would still
have to verify competitiveness against \emph{every} measurable model and
closedness of every comparison section.

\begin{thebibliography}{9}
\bibitem{PP} T.~Pennanen and A.-P.~Perkki\"o,
Topological duals of locally convex function spaces,
\emph{Positivity} 26, article 2 (2022), Theorem 1.
\href{https://doi.org/10.1007/s11117-022-00867-4}{doi:10.1007/s11117-022-00867-4}.
\end{thebibliography}
\end{document}
